\hypertarget{original-readme-start-here--what-every-file-in-this-folder-is-for}{%
\subsection{Start here --- what every file in this folder is
for}\label{original-readme-start-here--what-every-file-in-this-folder-is-for}}

This folder is meant to be opened as the project folder in Antigravity
(or any coding agent). Below is exactly what each file is, in the order
that matters, and what to actually do.

\hypertarget{original-readme-the-one-thing-to-understand-first}{%
\subsubsection{The one thing to understand
first}\label{original-readme-the-one-thing-to-understand-first}}

\texttt{hackstreak-spec.md}, \path|hackstreak-spec-v2.md|, and
\path|spec-v3-public-release.md| are \textbf{one design, revised three
times} --- not three products to build one after another. Each later
file states what it keeps and what it overrides from the one before it.
You will hand the agent all three at once; it reads them in order and
treats the later ones as corrections to the earlier ones.

What \emph{is} sequenced is build \textbf{phases} (Phase 1 → 1b → 2 → 3
→ 4), described inside the spec files themselves --- not the spec file
numbers. \texttt{kickoff-prompt.md} is written around this.

\textbf{This folder has been through an actual review pass} (reading
every file fresh and cross-checking them against each other, not just
re-reading what I remembered writing). It found and fixed five real
issues: a naming drift between \texttt{hack} (v1/v2) and \texttt{ascend}
(v3) that was stated once but repeated dozens of times the other way;
\texttt{hack done}/\texttt{hack mission ...} wording left over from
the old name in \path|new-unit-request-template.md|; a scope leak
where v2\textquotesingle s Phase 1b description includes a full Telegram
system that v3 actually narrows to desktop-only; a missing
\texttt{expanduser()} requirement that would have made the starter
pack\textquotesingle s own verifiers fail on night one, since its
\texttt{hack\_dir = "\textasciitilde{}/hack"} variable is substituted
as a literal string; and a lint rule (\texttt{pack.kind}) that was
described in prose but never added to the actual lint checklist. All
five are patched --- see "Implementation clarifications from final
review" near the top of \texttt{kickoff-prompt.md} for the
authoritative, consolidated fix list.

\textbf{One open item, left open on purpose:} the
project\textquotesingle s name. \path|spec-v3-public-release.md|
proposed "Ascend" as a placeholder and explicitly left the final choice
to you. Nothing has been built yet, so renaming later is a
find-and-replace, not a migration --- "Ascend" is used throughout as the
working name unless you say otherwise.

\hypertarget{original-readme-file-by-file}{%
\subsubsection{File-by-file}\label{original-readme-file-by-file}}

\begin{longtable}[]{@{}P{.24\textwidth}P{.35\textwidth}P{.33\textwidth}@{}}
\toprule\noalign{}
File & What it is & When it matters \\
\midrule\noalign{}
\endhead
\bottomrule\noalign{}
\endlastfoot
\textbf{\texttt{kickoff-prompt.md}} & The exact text to paste into
Antigravity to start. & Open this first. Copy its "Step 1" prompt block
and paste it as your first message to the agent, with this whole folder
open as the project. \\
\texttt{hackstreak-spec.md} & Original design: why this exists, the
streak math, the core CLI rules. Mostly superseded by the next file, but
still holds the streak algorithm and its test table, which stay valid. &
Read by the agent, not by you, unless you want the reasoning behind the
streak rules. \\
\path|hackstreak-spec-v2.md| & Adds: packs (the editable TOML roadmap
format), levels, milestones, verifiers (auto-checks). This is the core
engine design. Supersedes v1 wherever they disagree. & Read by the
agent. \\
\path|spec-v3-public-release.md| & Adds: making it
public/self-hostable, multi-user, the server + Docker architecture, and
an honest audit of what\textquotesingle s missing and
what\textquotesingle s deliberately deferred. Supersedes v2 on naming,
personal data, and deployment. & Read by the agent now (for context),
but its \emph{server} content (Phase 2+) isn\textquotesingle t built
until Step 3 below. \\
\texttt{ui-ux-spec.md} & The precise design brief for the web client:
screen layout, the level-path visual, colors, responsive behavior. & Not
used in Step 1 at all. Only matters once you reach Phase 2 (the web
client). Keep it in the folder anyway so the agent has full context. \\
\path|new-unit-request-template.md| & A template \emph{you} fill in
and paste into a \textbf{new chat with me (Claude)} whenever you want a
new learning pack (networking, DSA, a language, etc.) drafted. Not for
the coding agent. & Use later, whenever you want a new subject pack. Not
part of tonight\textquotesingle s build. \\
\path|packs/home-hack-lab.toml| & The one starter roadmap: build a
home hack lab on your Pi, attack it, harden it. A draft --- review it
yourself before trusting it (see its own header comments). & The agent
installs this as the sample pack during Step 1. \\
\texttt{docker/Dockerfile}, \path|docker/docker-compose.yml|,
\texttt{docker/.dockerignore} & How to run the server on your Pi (or
anywhere) once Phase 2 exists. & Not used until Step 3.
Don\textquotesingle t let the agent build \texttt{server/} before
then. \\
\end{longtable}

\hypertarget{original-readme-what-to-actually-do-in-order}{%
\subsubsection{What to actually do, in
order}\label{original-readme-what-to-actually-do-in-order}}

\textbf{Step 1 --- tonight.} Open \texttt{kickoff-prompt.md}, copy the
prompt under "Step 1," paste it to the agent with this folder as the
project. This builds \emph{only} the standalone CLI (Phase 1) and local
nudging (Phase 1b) --- no server, no web app yet.

\textbf{Step 2 --- before trusting it.} Run the tests yourself
(\texttt{pytest -v}), read
\path|packs/home-hack-lab.toml|\textquotesingle s \texttt{done\_when}
lines, run \texttt{ascend pack lint} and
\path|ascend pack check-links| on it, and actually use the CLI daily
for at least a week or two. This step is not optional --- a demo that
runs once in a chat is not the same as a habit tool that survives a
deadline week.

\textbf{Step 3 --- only once Step 2 has genuinely stuck.} Go back to the
agent with \path|spec-v3-public-release.md| and \texttt{ui-ux-spec.md}
and tell it to build Phase 2: the server (using the \texttt{docker/}
files) and the web client. Full instructions for this step are at the
end of \texttt{kickoff-prompt.md}.

\textbf{Later, whenever you want a new subject.} Fill in
\path|new-unit-request-template.md| and paste it into a \emph{new}
chat with me --- not the coding agent --- to get a new pack drafted,
then lint/review/install it the same way.

Don\textquotesingle t skip ahead. The whole point of staging this is
that a server and a phone app for a habit you haven\textquotesingle t
actually kept yet is building the sequel before the first movie was any
good.
